\chapter{باب كيف البيان وأقسامه}

\section{مقدمة الدرس وتعريف البيان والمجمل}
الحمد لله رب العالمين، والصلاة والسلام على رسول الله، أما بعد؛ فهذا هو الدرس الثاني من دروس شرح كتاب «الرسالة» للإمام أبي عبد الله الشافعي رحمه الله تعالى. ونشرع اليوم في «باب كيف البيان».

يقول الشافعي رحمه الله: «والبيان اسم جامع لمعانٍ مجتمعة الأصول متشعبة الفروع».
لفهم هذا الباب لابد من معرفة المجمل والمبين.
المجمل في اللغة: المجموع. وفي الاصطلاح: ما احتمل معنيين أو أكثر من غير ترجيح لواحد منهما.
أما البيان في اللغة: فهو الإيضاح والإظهار لما كان فيه خفاء. 

وقد اشترط بعض الأصوليين سبق الخفاء ليسمى بياناً، بينما اعتبر آخرون أن كل ما يزيل الإشكال ويوضح المعنى فهو بيان. والشافعي يرى أن البيان اسم يجمع تحته أصولاً مجتمعة تتشعب إلى فروع، يدركها ويفهمها أهل اللغة العربية، لأن القرآن نزل بلسانهم. فأقل ما في هذه المعاني أنها بيان لمن خوطب بها، متقاربة الاستواء في الفهم عندهم، وإن كان بعضها أشد تأكيداً في البيان من بعض. ومن يجهل لسان العرب قد يظن فيها الاختلاف والتشعب. وفي هذا حث لطلبة العلم على تعلم اللغة العربية والتبحر فيها لفهم مراد الله ورسوله.

\section{وجوه البيان عند الإمام الشافعي}
يقول الشافعي رحمه الله: «فجماع ما أبان الله لخلقه في كتابه مما تعبدهم به لما مضى من حكمه جل ثناؤه من وجوه». وذكر وجوه البيان، وهي إجمالاً:
\begin{enumerate}
    \item ما أبانه لخلقه نصاً في كتابه.
    \item ما أحكم فرضه بكتابه وبينه على لسان نبيه صلى الله عليه وسلم.
    \item ما سنه رسول الله صلى الله عليه وسلم مما ليس لله فيه نص حكم في القرآن.
    \item ما فرض الله على خلقه الاجتهاد في طلبه.
\end{enumerate}

\subsection{الاجتهاد في طلب الحق وإعمال النصوص}
قبل التفصيل في أقسام البيان، أشار الشافعي إلى النوع المتعلق بالاجتهاد. فقد ابتلانا الله بطاعته في الفرائض، وابتلانا كذلك بطاعته في الاجتهاد للوصول إلى الحق في بعض المسائل. ومن ذلك الاجتهاد في استقبال القبلة لمن غاب عن عين المسجد الحرام؛ قال تعالى: ﴿وَحَيْثُ مَا كُنتُمْ فَوَلُّوا وُجُوهَكُمْ شَطْرَهُ﴾. ففرض علينا الاجتهاد بالعقول التي ركبها فينا، والعلامات التي نصبها لنا كالشمس والقمر والنجوم والجبال والرياح، لنهتدي بها للقبلة. 

وكذلك الاجتهاد في معرفة «العدل» لقبول شهادته أو حكمه؛ قال تعالى: ﴿وَأَشْهِدُوا ذَوَيْ عَدْلٍ مِّنكُمْ﴾، والعدل هو المسلم البالغ العاقل الذي سلم من أسباب الفسق وخوارم المروءة، فنجتهد في تطبيق هذه الشروط على الأشخاص.
وكذلك الاجتهاد في «جزاء الصيد»؛ كأن يصيد المحرم صيداً، فيجتهد الحكمان العدلان في تحديد المثل من النعم الذي يقابل هذا الصيد ليكون فدية. وهذا النوع من الاجتهاد يسميه الشافعي «القياس»، وهو القياس الصحيح الذي يعتمد على إعمال النصوص والقرائن، بخلاف القول بالاستحسان المجرد عن الدليل. ومن هذا يتبين أنه لا يجوز لأحد (غير رسول الله صلى الله عليه وسلم) أن يحلل أو يحرم إلا بالاستدلال من الكتاب أو السنة أو الإجماع أو القياس الصحيح.

\section{تفصيل أبواب البيان}

\subsection{البيان الأول: ما أبانه الله لخلقه نصاً في كتابه}
وهو ما جاء نصاً واضحاً لا يحتمل تأويلاً، كفرضية الصيام، والصلاة، والزكاة، والحج، وكتحريم الزنا والخمر والميتة والدم. وضرب الشافعي لذلك أمثلة تتضمن زيادة في البيان اللفظي والتأكيد، كقوله تعالى: ﴿فَمَن لَّمْ يَجِدْ فَصِيَامُ ثَلَاثَةِ أَيَّامٍ فِي الْحَجِّ وَسَبْعَةٍ إِذَا رَجَعْتُمْ ۗ تِلْكَ عَشَرَةٌ كَامِلَةٌ﴾، وقوله: ﴿وَوَاعَدْنَا مُوسَىٰ ثَلَاثِينَ لَيْلَةً وَأَتْمَمْنَاهَا بِعَشْرٍ فَتَمَّ مِيقَاتُ رَبِّهِ أَرْبَعِينَ لَيْلَةً﴾. فهذا من باب زيادة التبيين لجملة العدد لمن خوطبوا بالآية.

\subsection{البيان الثاني: ما نص القرآن على فرضه وبينت السنة بعض أحكامه}
وهو ما أتى فرضه في القرآن، ولكن بعض جزئياته أو شروطه احتاجت إلى تقييد أو تخصيص أو بيان من السنة. وضرب له أمثلة:
\begin{itemize}
    \item \textbf{الوضوء والغسل:} قال تعالى: ﴿إِذَا قُمْتُمْ إِلَى الصَّلَاةِ فَاغْسِلُوا وُجُوهَكُمْ وَأَيْدِيَكُمْ إِلَى الْمَرَافِقِ﴾ الآية. فأتى القرآن بالفرض، لكن السنة بينت الكيفية، وأن العدد يجزئ فيه المرة، والكمال في الثلاث، وبينت موجبات الوضوء والغسل، ووضحت أن غسل الأعقاب واجب في الوضوء بقوله صلى الله عليه وسلم: «ويل للأعقاب من النار».
    \item \textbf{المواريث والوصية:} فصل القرآن أنصبة المواريث تفصيلاً دقيقاً، لكن اشترط أن يكون التقسيم ﴿مِن بَعْدِ وَصِيَّةٍ يُوصِي بِهَا أَوْ دَيْنٍ﴾. فجاءت السنة المطهرة في حديث سعد بن أبي وقاص رضي الله عنه لتقيد الوصية بألا تتجاوز الثلث، فقال صلى الله عليه وسلم: «الثلث، والثلث كثير».
\end{itemize}

\subsection{البيان الثالث: ما أحكم الله فرضه مجملاً وبينت السنة كيفيته}
وهذا النوع يختلف عن سابقه في أن الفرض يأتي في القرآن مجملاً إجمالاً تاماً، وتتولى السنة تفصيل كل أحكامه وكيفياته. كقوله تعالى: ﴿إِنَّ الصَّلَاةَ كَانَتْ عَلَى الْمُؤْمِنِينَ كِتَابًا مَّوْقُوتًا﴾، وقوله: ﴿وَأَقِيمُوا الصَّلَاةَ وَآتُوا الزَّكَاةَ﴾، و﴿وَلِلَّهِ عَلَى النَّاسِ حِجُّ الْبَيْتِ﴾. فبين النبي صلى الله عليه وسلم عدد الصلوات ومواقيتها وسننها بقوله: «صلوا كما رأيتموني أصلي»، وبين أنصبة الزكاة، وبين مناسك الحج بقوله: «خذوا عني مناسككم». 

\subsection{البيان الرابع: ما سنه النبي صلى الله عليه وسلم استقلالاً}
وهو كل ما سنه رسول الله صلى الله عليه وسلم مما ليس فيه نص كتاب مباشر. كتحريم الجمع بين المرأة وعمتها، والمرأة وخالتها، وغيرها من التشريعات. وهذا البيان حجة وفرض يجب قبوله، لأن الله تعالى فرض في كتابه طاعة رسوله طاعة مطلقة، قال تعالى: ﴿وَمَا آتَاكُمُ الرَّسُولُ فَخُذُوهُ وَمَا نَهَاكُمْ عَنْهُ فَانتَهُوا﴾. 

فمن ادعى أنه يتبع القرآن فقط وترك السنة، فقد كذب القرآن الذي أمره باتباع السنة. ومن قَبِلَ عن رسول الله صلى الله عليه وسلم سنته، فقد قَبِلَ عن الله عز وجل أمره وفرضه. ومن فرّق بينهما، فقد ردّ على الله الذي فرض طاعة نبيه. فكل شيء شرعه النبي صلى الله عليه وسلم فهو بيان من الحكمة التي آتاه الله إياها، ويجب التسليم له والتصديق به.